\apartado{Fundamento cloud y modelo de servicio}
\subsection{SaaS para el estudiante; plataforma administrada para el equipo}
El NIST caracteriza la nube mediante acceso por red, recursos compartidos, elasticidad y medición del uso, entre otros rasgos. Distingue SaaS, PaaS e IaaS según la capacidad entregada y la responsabilidad del consumidor \cite{nist}. En el proyecto, el estudiante consumiría una \textbf{aplicación como servicio (SaaS)}: accedería por web a sus materias y datos sin administrar servidores. Esta clasificación describe la experiencia prevista, no una certificación ni un servicio ya publicado.

Para construir esa aplicación se propone utilizar capacidades de plataforma y servicios administrados. Cloud Run ejecuta contenedores y administra instancias; Firestore, Storage y Firebase resuelven persistencia e identidad. El equipo conserva la responsabilidad de sus contratos, permisos, dependencias, reglas de negocio y configuración. Usar un servicio administrado no delega la decisión de quién puede leer una materia.

\begin{tabularx}{\linewidth}{P{2.6cm}Y Y}
\toprule
\textbf{Capa} & \textbf{Necesidad del proyecto} & \textbf{Responsabilidad del equipo}\\
\midrule
Aplicación SaaS & Materias y agenda privadas disponibles por navegador. & Experiencia, accesibilidad, propiedad y consistencia.\\
Plataforma de contenedores & Ejecutar web, API y worker sin mantener una VM propia. & Imagen, dependencias, límites, entradas y versiones.\\
Datos administrados & Conservar registros y archivos más allá de una instancia. & Modelo, autorización, índices, retención y borrado.\\
Identidad e IAM & Identificar personas y autenticar servicios. & Roles mínimos, cuentas del piloto y pruebas negativas.\\
Operación administrada & Observar uso, fallos y consumo. & Señales útiles, redacción, alertas y respuesta.\\
\bottomrule
\end{tabularx}

\subsection{Por qué esta arquitectura}
El monolito modular reduce el número de contratos distribuidos durante el piloto. API y worker comparten dominio, pero pueden tener despliegue y capacidad distintos: navegar no debe competir con una extracción pesada. Se evita dividir cada función en un microservicio antes de comprobar que lo necesita.

La propuesta usa una nube pública para alojar un piloto de acceso restringido. \textbf{Piloto privado no equivale a nube privada}: la privacidad de los datos depende de autenticación, autorización y configuración; no del nombre dado al entorno. El entorno local facilita desarrollo y pruebas, mientras que el entorno cloud futuro permite comprobar identidades de servicio, persistencia y elasticidad reales.

\destacado{La nube como parte del sistema}{El aporte previsto incluye persistencia multiusuario, acceso web, separación de cargas, procesamiento asíncrono, medición y trazabilidad. Alojar una maqueta estática no demostraría por sí solo estas capacidades.}
